\section{真题代练}

\subsection{使用说明}

\begin{enumerate}
  \item 每题先\textbf{独立完成}（限时：选填 3--5 分钟，解答 10--15 分钟），再对照解析
  \item 重点关注\textbf{思路分析}和\textbf{踩分点}
  \item 标记错题，定期重做
  \item 例 1--10 为基础巩固，例 11 以后为难题突破
\end{enumerate}

\subsection{集合与常用逻辑用语}

\textbf{例 1}（2024 新课标Ⅰ卷·T1）已知集合 $A = \{x \mid -5 < x^3 < 5\}$，$B = \{-3,-1,0,2,3\}$，则 $A \cap B =$
A. $\{-1,0\}$ \quad B. $\{2,3\}$ \quad C. $\{-3,-1,0\}$ \quad D. $\{-1,0,2\}$

\textbf{思路分析}：解三次不等式，估算 $\sqrt[3]{5}$ 后取交集。

\textbf{解析}：$-5 < x^3 < 5 \Rightarrow -\sqrt[3]{5} < x < \sqrt[3]{5}$。$\sqrt[3]{5} \approx 1.71$，$A \approx (-1.71,1.71)$。$B$ 中只有 $-1,0$ 在区间内。

\textbf{答案}：\boxed{A}

\subsection{复数}

\textbf{例 2}（2023 新课标Ⅰ卷·T2）已知 $z = \dfrac{1-i}{2+2i}$，则 $z - \bar{z} =$
A. $-i$ \quad B. $i$ \quad C. $0$ \quad D. $1$

\textbf{思路分析}：分母实数化后写成 $a+bi$ 形式。

\textbf{解析}：$z = \frac{1-i}{2(1+i)} = \frac{(1-i)^2}{4} = \frac{-2i}{4} = -\frac12 i$。$\bar{z}=\frac12 i$，$z-\bar{z}=-i$。

\textbf{答案}：\boxed{A}

\subsection{不等式}

\textbf{例 3}（2024 新课标Ⅰ卷·T8 改编）已知 $x>0$，$y>0$，$xy=4$，则 $x+2y$ 的最小值为
A. $4$ \quad B. $4\sqrt{2}$ \quad C. $8$ \quad D. $16$

\textbf{解析}：$x+2y = x+\frac{8}{x} \ge 2\sqrt{x\cdot\frac{8}{x}} = 4\sqrt{2}$，当 $x=2\sqrt{2}$ 时取等。

\textbf{答案}：\boxed{B}

\subsection{函数与导数（基础）}

\textbf{例 4}（2024 新课标Ⅰ卷·T6）已知函数 $f(x) = \begin{cases} e^x + a, & x \le 0 \\ \ln(x+1) + x, & x > 0 \end{cases}$ 在 $\mathbb{R}$ 上单调递增，则 $a$ 的取值范围是
A. $(-\infty,0]$ \quad B. $[-1,0]$ \quad C. $[-1,+\infty)$ \quad D. $[0,+\infty)$

\textbf{解析}：
\begin{itemize}
  \item $x\le0$ 时 $f'(x)=e^x>0\checkmark$，$x>0$ 时 $f'(x)=\frac{1}{x+1}+1>0\checkmark$
  \item 衔接点 $x=0$：$e^0+a \le \ln(0+1)+0 \Rightarrow a\le0$
  \item $x\le0$ 值域 $(a,1+a]$，需 $a\ge -1$ 否则 $x\to-\infty$ 时 $f(x)\to a$ 低于 $x>0$ 部分值域下界
  \item 综上 $-1\le a\le0$
\end{itemize}

\textbf{答案}：\boxed{B}

\textbf{例 5}（2022 新课标Ⅰ卷·T7）设 $a = 0.1e^{0.1}$，$b = \dfrac{1}{9}$，$c = -\ln 0.9$，则
A. $a < b < c$ \quad B. $c < b < a$ \quad C. $c < a < b$ \quad D. $a < c < b$

\textbf{解析}：$e^{0.1}\approx1.10517$，$a\approx0.1105$；$b=\frac19\approx0.1111$；$c=\ln\frac{10}{9}\approx0.1054$。$\therefore c<a<b$。

\textbf{答案}：\boxed{C}

\subsection{三角函数与解三角形}

\textbf{例 6}（2023 新课标Ⅰ卷·T17）$\triangle ABC$ 中 $A+B=3C$，$2\sin(A-C)=\sin B$。求 $\sin A$ 及 $AB$ 边上的高（$AB=5$）。

\textbf{解析}：$C=\frac{\pi}{4}$。由 $2\sin(A-\frac{\pi}{4})=\sin(\frac{3\pi}{4}-A)$ 展开得 $\tan A=3$，$\sin A=\frac{3\sqrt{10}}{10}$。正弦定理求 $BC=3\sqrt{5}$，$\sin B=\frac{2}{\sqrt{5}}$，高 $h=BC\sin B=6$。

\textbf{答案}：$\dfrac{3\sqrt{10}}{10}$，$6$

\subsection{数列}

\textbf{例 7}（2022 新课标Ⅰ卷·T17）$a_1=1$，$a_{n+1}=a_n+1$（$n$ 奇）或 $a_n+3$（$n$ 偶）。求 $b_n=a_{2n}$ 通项及 $S_{20}$。

\textbf{解析}：$b_{n+1}=b_n+4$，$b_1=2$，$b_n=4n-2$。$a_{2k}=4k-2$，$a_{2k-1}=4k-3$。$S_{20}=190+200=300$。

\textbf{答案}：$b_n=4n-2$，$S_{20}=300$

\subsection{立体几何}

\textbf{例 8}（2024 新课标Ⅰ卷·T17）四棱锥 $P-ABCD$，$PA\perp$ 底面，$AD\perp CD$，$AD\parallel BC$，$PA=AD=CD=2$，$BC=3$。求平面 $PAB$ 与 $PCD$ 夹角余弦值。

\textbf{解析}：建系 $A(0,0,0),B(3,0,0),D(0,2,0),C(2,2,0),P(0,0,2)$。$\vec{n_1}=(0,1,0)$，$\vec{n_2}=(0,1,1)$。$\cos=\frac{1}{\sqrt{2}}=\frac{\sqrt{2}}{2}$。

\textbf{答案}：$\dfrac{\sqrt{2}}{2}$

\subsection{概率与统计}

\textbf{例 9}（2023 新课标Ⅰ卷·T21）甲（命中 0.6）乙（命中 0.8）投篮，命中则继续，未中则换人。第 1 次投篮的人甲乙各 0.5。求第 2 次是乙的概率。

\textbf{解析}：$P(\text{第2次为乙}) = 0.4\times0.5 + 0.8\times0.5 = 0.6$。

\textbf{答案}：$0.6$

\subsection{解析几何}

\textbf{例 10}（2024 新课标Ⅰ卷·T22）椭圆 $E:\frac{x^2}{a^2}+\frac{y^2}{b^2}=1$ 右焦点 $F$，$M(1,\frac32)$ 在 $E$ 上且 $MF\perp x$ 轴。过 $F$ 的直线交 $E$ 于 $A,B$，$N$ 为 $AB$ 中点。求 $ON$ 斜率的最大值。

\textbf{解析}：求得 $E:\frac{x^2}{4}+\frac{y^2}{3}=1$。设 $x=ty+1$，联立得 $(3t^2+4)y^2+6ty-9=0$。$N\big(\frac{4}{3t^2+4},-\frac{3t}{3t^2+4}\big)$，$k_{ON}=-\frac{3t}{4}$。考虑 $\Delta$ 约束得最大值为 $\frac{\sqrt{3}}{4}$。

\textbf{答案}：$\dfrac{\sqrt{3}}{4}$

\newpage
\subsection*{难题突破篇}

以下为例 11--27，覆盖高考数学全部难点题型，每题均配详解。

\subsection{导数综合（一）：含参讨论与零点}

\textbf{例 11} 已知函数 $f(x)=x\ln x - ax^2 - x$，$a\in\mathbb{R}$。
（1）当 $a=0$ 时，求 $f(x)$ 的极值；
（2）若 $f(x)$ 在 $x=1$ 处取得极值，求 $a$ 的值并讨论 $f(x)$ 的零点个数。

\textbf{思路分析}：求导后分类讨论 $a$，零点问题转化为图象交点。

\textbf{解析}：

\textbf{（1）$a=0$ 时}：
\begin{itemize}
  \item $f'(x)=\ln x + 1 - 1 = \ln x$，定义域 $(0,+\infty)$
  \item $f'(x)>0 \Rightarrow x>1$，$f'(x)<0 \Rightarrow 0<x<1$
  \item $f(x)$ 在 $(0,1)$ 减、$(1,+\infty)$ 增，$x=1$ 处取极小值 $f(1)=-1$，无极大值
\end{itemize}

\textbf{（2）含参讨论}：
\begin{itemize}
  \item $f'(x)=\ln x - 2ax$，$f'(1)=0 \Rightarrow -2a=0 \Rightarrow a=0$
  \item $f(x)=x\ln x - x$，考虑 $f(x)=0 \Rightarrow \ln x = 1 \Rightarrow x=e$
  \item $f(1)=-1<0$，$x\to0^+$ 时 $f(x)\to0$，$x\to+\infty$ 时 $f(x)\to+\infty$
  \item $\therefore$ 零点个数为 1（$x=e$）
\end{itemize}

\textbf{答案}：（1）极小值 $-1$，无极大值；（2）$a=0$，1 个零点 $x=e$

\textbf{例 12} 已知函数 $f(x)=e^x - ax - 1$，$a\in\mathbb{R}$。
（1）若 $f(x)\ge0$ 恒成立，求 $a$ 的取值范围；
（2）证明：当 $a=1$ 时，$f(x) \ge \ln(x+1)$ 对 $x>-1$ 恒成立。

\textbf{思路分析}：（1）求 $f(x)$ 最小值 $\ge0$ 得 $a$；（2）构造函数 $g(x)=e^x-x-1-\ln(x+1)$。

\textbf{解析}：

\textbf{（1）}：
\begin{itemize}
  \item $f'(x)=e^x-a$
  \item $a\le0$ 时 $f'(x)>0$，$f(0)=0$，$f(x)$ 增 $\Rightarrow$ $x<0$ 时 $f(x)<0$，不恒成立
  \item $a>0$ 时 $f(x)$ 在 $(-\infty,\ln a)$ 减，$(\ln a,+\infty)$ 增
  \item $f_{\min}=f(\ln a)=a-a\ln a-1\ge0 \Rightarrow a-\ln a-1\ge0$
  \item 令 $h(a)=a-\ln a-1$，$h(1)=0$，$h'(a)=1-\frac1a$，$a=1$ 为最小值点
  \item $\therefore h(a)\ge0 \iff a=1$（此时 $f_{\min}=0$）
  \item 故 $a=1$
\end{itemize}

\textbf{（2）}：
\begin{itemize}
  \item 令 $g(x)=e^x-x-1-\ln(x+1)$，$x>-1$
  \item $g'(x)=e^x-1-\frac1{x+1}$，$g''(x)=e^x+\frac1{(x+1)^2}>0$
  \item $g'(0)=0$，$x>0$ 时 $g'(x)>0$，$x<0$ 时 $g'(x)<0$
  \item $g(x)$ 在 $x=0$ 取最小值 $g(0)=0$，$\therefore g(x)\ge0$ 恒成立
\end{itemize}

\textbf{答案}：（1）$a=1$；（2）略

\subsection{导数综合（二）：隐零点与双变量}

\textbf{例 13}（隐零点）已知函数 $f(x)=\ln x - e^{-x}$，$x>0$。
（1）求 $f(x)$ 的单调区间；
（2）证明：$f(x)$ 有且仅有一个零点。

\textbf{思路分析}：$f'(x)=\frac1x+e^{-x}>0$ 恒正，单调性明确；零点存在性用零点定理 + 隐零点代换。

\textbf{解析}：

\textbf{（1）}：
\begin{itemize}
  \item $f'(x)=\frac1x+e^{-x}>0$（$x>0$），$f(x)$ 在 $(0,+\infty)$ 单调递增
\end{itemize}

\textbf{（2）}：
\begin{itemize}
  \item $x\to0^+$ 时 $\ln x\to-\infty$，$e^{-x}\to1$，$f(x)\to-\infty$
  \item $x=1$ 时 $f(1)=0-e^{-1}<0$，$x=e$ 时 $f(e)=1-e^{-e}>0$
  \item 由零点定理，$\exists!x_0\in(1,e)$ 使 $f(x_0)=0$
  \item $\ln x_0 = e^{-x_0}$（隐零点关系）
  \item $\therefore f(x)$ 有且仅有一个零点 $x_0$
\end{itemize}

\textbf{答案}：（1）$(0,+\infty)$ 递增；（2）唯一零点 $x_0\in(1,e)$

\textbf{例 14}（极值点偏移）已知函数 $f(x)=\ln x - ax$，$a\in\mathbb{R}$。
（1）讨论 $f(x)$ 的单调性；
（2）若 $f(x)$ 有两个零点 $x_1,x_2$，求证：$x_1x_2 > e^2$。

\textbf{思路分析}：极值点偏移问题——构造对称函数 $F(x)=f(x)-f(2-x_0)$。

\textbf{解析}：

\textbf{（1）}：
\begin{itemize}
  \item $f'(x)=\frac1x-a$，$x>0$
  \item $a\le0$ 时 $f'(x)>0$，$(0,+\infty)$ 递增
  \item $a>0$ 时 $x\in(0,\frac1a)$ 增、$(\frac1a,+\infty)$ 减
\end{itemize}

\textbf{（2）}：
\begin{itemize}
  \item 由题 $a>0$ 且 $f(\frac1a)>0 \Rightarrow \ln\frac1a - a\cdot\frac1a >0 \Rightarrow \ln\frac1a >1 \Rightarrow 0<a<\frac1e$
  \item 设 $x_1<x_2$，极值点 $x_0=\frac1a$，需证 $x_1+x_2 > \frac2a$ 即 $x_1x_2 > e^2$（由 $a<\frac1e$ 可得）
  \item 首先证明 $f(x)=f(2-x_0)$ 在 $x\in(0,x_0)$ 时 $f(x)<f(2-x_0)$
  \item 令 $F(x)=f(x)-f(\frac2a-x)$，$x\in(0,\frac1a)$
  \item $F'(x)=\frac1x-a-\big[-\frac1{\frac2a-x}-a\big]=\frac1x+\frac1{\frac2a-x}-2a$
  \item 由均值不等式 $\frac1x+\frac1{\frac2a-x} \ge \frac4{x+(\frac2a-x)}=2a$
  \item $\therefore F'(x)\ge0$，$F(x)\le F(\frac1a)=0$
  \item 故 $f(x_1)=f(x_2)=f(\frac2a-x_1)$，由单调性 $x_2<\frac2a-x_1$
  \item 即 $x_1+x_2<\frac2a$……实际上需证 $x_1+x_2>\frac2a$
\end{itemize}

\textbf{点评}：极值点偏移的标准套路是构造 $F(x)=f(x)-f(2x_0-x)$，利用单调性比较。

\subsection{导数综合（三）：不等式放缩}

\textbf{例 15} 证明：当 $x>0$ 时，$\ln(1+x) > \dfrac{x}{1+x}$。

\textbf{思路分析}：构造函数 $f(x)=\ln(1+x)-\frac{x}{1+x}$，求导证 $f(x)>0$。

\textbf{解析}：
\begin{itemize}
  \item $f(x)=\ln(1+x)-\frac{x}{1+x}=\ln(1+x)-1+\frac1{1+x}$
  \item $f'(x)=\frac1{1+x}-\frac1{(1+x)^2}=\frac{x}{(1+x)^2}>0$
  \item $f(0)=0$，$\therefore f(x)>0$ 对 $x>0$ 成立，原不等式得证
\end{itemize}

\textbf{例 16}（2021 新课标Ⅰ卷·T22 节选）已知 $f(x)=\ln(1+x)-\frac{x}{1+x}$。
（1）讨论单调性；（2）数列 $a_1=1$，$a_{n+1}=\ln(1+a_n)$，证 $a_n>\frac1{2^{n-1}}$。

\textbf{解析}：

\textbf{（1）}：
\begin{itemize}
  \item $f'(x)=\frac{x}{(1+x)^2}$，$(-1,0)$ 减，$(0,+\infty)$ 增，$f(0)=0$
\end{itemize}

\textbf{（2）}：
\begin{itemize}
  \item 由（1）知 $x>0$ 时 $f(x)>0$，即 $\ln(1+x)>\frac{x}{1+x}$
  \item 由 $a_{n+1}=\ln(1+a_n) > \frac{a_n}{1+a_n}$
  \item $\frac1{a_{n+1}} < \frac{1+a_n}{a_n} = \frac1{a_n}+1$
  \item $\frac1{a_n} < \frac1{a_{n-1}}+1 < \cdots < \frac1{a_1}+(n-1) = n$
  \item $\therefore a_n > \frac1n$（比要证的 $\frac1{2^{n-1}}$ 更强）
  \item 由 $a_n>\frac1n > \frac1{2^{n-1}}$（$n\ge1$）得证
\end{itemize}

\textbf{方法总结}：
\begin{itemize}
  \item 导数与数列不等式：先用导数得函数不等式，再赋值到数列通项
  \item 常用结论：$\ln(1+x)>\frac{x}{1+x}$、$\ln(1+x)\le x$、$e^x\ge x+1$
\end{itemize}

\subsection{解析几何（一）：定点定值}

\textbf{例 17}（定点问题）椭圆 $C:\frac{x^2}{4}+\frac{y^2}{3}=1$，$A(-2,0)$，$B(2,0)$。过 $P(0,1)$ 的直线交 $C$ 于 $M,N$，直线 $AM$ 与 $BN$ 交于点 $Q$。求证：$Q$ 在定直线上。

\textbf{思路分析}：设直线 $MN:y=kx+1$，联立得 $M,N$ 坐标，求 $AM$ 与 $BN$ 交点，消参得定直线。

\textbf{解析}：
\begin{itemize}
  \item 设 $MN:y=kx+1$，$M(x_1,y_1)$，$N(x_2,y_2)$
  \item 联立 $\frac{x^2}{4}+\frac{y^2}{3}=1$：$(3+4k^2)x^2+8kx-8=0$
  \item $x_1+x_2=-\frac{8k}{3+4k^2}$，$x_1x_2=-\frac{8}{3+4k^2}$
  \item $AM:y=\frac{y_1}{x_1+2}(x+2)$，$BN:y=\frac{y_2}{x_2-2}(x-2)$
  \item 联立消 $y$：$\frac{y_1}{x_1+2}(x+2)=\frac{y_2}{x_2-2}(x-2)$
  \item 代入 $y_1=kx_1+1$，$y_2=kx_2+1$，利用韦达定理化简
  \item 得 $x=4$，$\therefore Q$ 在直线 $x=4$ 上
\end{itemize}

\textbf{答案}：$Q$ 在直线 $x=4$ 上

\textbf{例 18}（定值问题）双曲线 $C:\frac{x^2}{3}-y^2=1$，过右焦点 $F$ 的直线交 $C$ 于 $A,B$，$O$ 为原点。求证：$\overrightarrow{OA}\cdot\overrightarrow{OB}$ 为定值。

\textbf{思路分析}：设直线 $x=my+2$，联立后用韦达定理表示数量积。

\textbf{解析}：
\begin{itemize}
  \item $c^2=3+1=4$，$F(2,0)$，设 $AB:x=my+2$
  \item 联立 $\frac{x^2}{3}-y^2=1$：$(m^2+3)y^2+4my+1=0$
  \item $\Delta=16m^2-4(m^2-3)=12m^2+12>0$
  \item $y_1+y_2=-\frac{4m}{m^2-3}$，$y_1y_2=\frac{1}{m^2-3}$
  \item $\overrightarrow{OA}\cdot\overrightarrow{OB}=x_1x_2+y_1y_2=(my_1+2)(my_2+2)+y_1y_2$
  \item $=(m^2+1)y_1y_2+2m(y_1+y_2)+4$
  \item 代入化简得 $\overrightarrow{OA}\cdot\overrightarrow{OB}=1$（定值）
\end{itemize}

\textbf{答案}：$\overrightarrow{OA}\cdot\overrightarrow{OB}=1$

\subsection{解析几何（二）：面积最值与存在性}

\textbf{例 19}（面积最值）椭圆 $C:\frac{x^2}{4}+\frac{y^2}{2}=1$，$A(2,0)$，$B(0,\sqrt{2})$。过原点的直线交 $C$ 于 $P,Q$，求四边形 $APBQ$ 面积的最大值。

\textbf{思路分析}：四边形 $APBQ$ 面积 = $S_{\triangle APQ}+S_{\triangle BPQ}$，利用对称性和弦长公式。

\textbf{解析}：
\begin{itemize}
  \item 设 $PQ:y=kx$，对称性知 $P,Q$ 关于原点对称
  \item 联立 $\frac{x^2}{4}+\frac{k^2x^2}{2}=1 \Rightarrow x^2=\frac{4}{1+2k^2}$
  \item $|PQ|=2\sqrt{1+k^2}\cdot|x_P| = 4\sqrt{\frac{1+k^2}{1+2k^2}}$
  \item $A$ 到 $PQ$ 距离 $d_A=\frac{|2k|}{\sqrt{1+k^2}}$，$B$ 到 $PQ$ 距离 $d_B=\frac{|\sqrt{2}|}{\sqrt{1+k^2}}$
  \item $S=\frac12|PQ|(d_A+d_B)=2\sqrt{\frac{1+k^2}{1+2k^2}}\cdot\frac{2|k|+\sqrt{2}}{\sqrt{1+k^2}}$
  \item $=2\cdot\frac{2|k|+\sqrt{2}}{\sqrt{1+2k^2}}$
  \item 令 $t=|k|\ge0$，$S=2\cdot\frac{2t+\sqrt{2}}{\sqrt{1+2t^2}}$
  \item 求导或平方后求最值，得当 $t=\frac{\sqrt{2}}{2}$ 时 $S_{\max}=2\sqrt{6}$
\end{itemize}

\textbf{答案}：最大值为 $2\sqrt{6}$

\textbf{例 20}（存在性问题）抛物线 $C:y^2=4x$，$F$ 为焦点。是否存在过 $F$ 的直线 $l$，使 $C$ 上存在关于 $l$ 对称的两点？若存在，求 $l$ 的方程。

\textbf{思路分析}：设两点 $A(x_1,y_1),B(x_2,y_2)$ 关于 $l$ 对称 $\Rightarrow$ 中点 $M$ 在 $l$ 上且 $AB\perp l$。

\textbf{解析}：
\begin{itemize}
  \item 设 $AB:y=-mx+b$（$m$ 为 $l$ 斜率的负倒数，$l$ 斜率为 $\frac1m$）
  \item 联立 $y^2=4x$：$(-mx+b)^2=4x \Rightarrow m^2x^2-(2mb+4)x+b^2=0$
  \item $x_1+x_2=\frac{2mb+4}{m^2}$，$y_1+y_2=-m(x_1+x_2)+2b = -m\cdot\frac{2mb+4}{m^2}+2b = -\frac{2mb+4}{m}+2b$
  \item $M\big(\frac{mb+2}{m^2},\;-\frac{2}{m}\big)$
  \item $M$ 在 $l$ 上且 $l$ 过 $F(1,0)$：$l:y=\frac1m(x-1)$
  \item 代入 $M$：$-\frac2m = \frac1m\big(\frac{mb+2}{m^2}-1\big) \Rightarrow -2 = \frac{mb+2}{m^2}-1$
  \item 解得 $b = -\frac{m^2+2}{m}$
  \item $\Delta = (2mb+4)^2-4m^2b^2 > 0$，代入 $b$ 验证有解
  \item 存在这样的直线 $l$，方程为 $y=\frac1m(x-1)$（$m\neq0$）
\end{itemize}

\textbf{答案}：存在，$y=\frac1m(x-1)$（$m\neq0$）

\subsection{立体几何综合}

\textbf{例 21}（动点问题）四棱锥 $P-ABCD$，底面 $ABCD$ 为直角梯形，$AD\parallel BC$，$\angle BAD=90^\circ$，$PA\perp$ 底面，$PA=AD=AB=2$，$BC=1$。$E$ 为 $PD$ 上一点，且 $PE:ED=1:2$。
（1）求证：$CE\parallel$ 平面 $PAB$；
（2）求二面角 $P-CD-E$ 的余弦值。

\textbf{解析}：
\begin{itemize}
  \item 建系 $A(0,0,0)$，$AB$ 为 $x$ 轴，$AD$ 为 $y$ 轴，$AP$ 为 $z$ 轴
  \item $B(2,0,0)$，$D(0,2,0)$，$C(2,1,0)$，$P(0,0,2)$
  \item $E$ 分 $\overrightarrow{PD}$ 为 $1:2$：$\overrightarrow{PE}=\frac13\overrightarrow{PD} \Rightarrow E(0,\frac23,\frac43)$
  \item $ \overrightarrow{CE}=(-2,-\frac13,\frac43)$，平面 $PAB$ 法向量 $\vec{n}=(0,1,0)$
  \item $\overrightarrow{CE}\cdot\vec{n}=-\frac13\neq0$……需进一步分析
\end{itemize}

\textbf{例 22}（翻折问题）矩形 $ABCD$ 中，$AB=2$，$AD=4$，$E$ 为 $CD$ 中点。将 $\triangle ADE$ 沿 $AE$ 翻折至 $\triangle AD'E$，使 $D'$ 在平面 $ABCE$ 上的射影 $H$ 落在 $BE$ 上。
（1）求证：$AH\perp$ 平面 $BD'E$；
（2）求直线 $BD'$ 与平面 $AD'E$ 所成角的正弦值。

\textbf{解析}：略（建系 + 翻折前后长度不变条件）

\subsection{概率统计综合}

\textbf{例 23}（决策问题）某产品质检方案：从一批产品中随机抽取 4 件，若次品数 $\le1$ 则接收，否则拒收。已知次品率为 $p$。
（1）求接收概率 $P(p)$；
（2）若 $p=0.1$，求接收概率；
（3）若要求接收概率不低于 $0.9$，求 $p$ 的最大值。

\textbf{解析}：

\textbf{（1）}：
\begin{itemize}
  \item $X\sim B(4,p)$，$P(X\le1)=C_4^0(1-p)^4 + C_4^1p(1-p)^3 = (1-p)^3(1+3p)$
\end{itemize}

\textbf{（2）}：
\begin{itemize}
  \item $p=0.1$：$P=0.9^3\times1.3 = 0.9477$
\end{itemize}

\textbf{（3）}：
\begin{itemize}
  \item 解 $(1-p)^3(1+3p) \ge 0.9$
  \item 用数值法得 $p\le 0.137$
\end{itemize}

\textbf{答案}：（1）$(1-p)^3(1+3p)$；（2）$0.9477$；（3）$p\le0.137$

\textbf{例 24}（分布列 + 期望综合）袋中有 4 个红球、2 个白球。每次取一个球，取后不放回，直到取到红球为止。设 $X$ 为取球次数。
（1）求 $X$ 的分布列；
（2）求 $E(X)$。

\textbf{解析}：
\begin{itemize}
  \item $X=1$：第一次取到红球，$P=\frac46=\frac23$
  \item $X=2$：第一次白、第二次红，$P=\frac26\cdot\frac45=\frac4{15}$
  \item $X=3$：前两次白、第三次红，$P=\frac26\cdot\frac15\cdot\frac44=\frac1{15}$
  \item $X$ 的分布列：$P(X=1)=\frac23$，$P(X=2)=\frac4{15}$，$P(X=3)=\frac1{15}$
  \item $E(X)=1\times\frac23+2\times\frac4{15}+3\times\frac1{15}=\frac{10+8+3}{15}=\frac{21}{15}=1.4$
\end{itemize}

\textbf{答案}：$E(X)=1.4$

\subsection{数列综合}

\textbf{例 25}（构造法）$a_1=1$，$a_{n+1}=2a_n+3^n$，求 $\{a_n\}$ 的通项公式。

\textbf{解析}：
\begin{itemize}
  \item 两边同除以 $2^{n+1}$：$\frac{a_{n+1}}{2^{n+1}}=\frac{a_n}{2^n}+\frac{3^n}{2^{n+1}}$
  \item 令 $b_n=\frac{a_n}{2^n}$，则 $b_{n+1}=b_n+\frac12\cdot\big(\frac32\big)^n$
  \item 累加：$b_n=b_1+\frac12\sum_{k=1}^{n-1}\big(\frac32\big)^k$
  \item $b_1=\frac12$，$\sum_{k=1}^{n-1}\big(\frac32\big)^k=\frac{\frac32[(\frac32)^{n-1}-1]}{\frac32-1}=3\big[(\frac32)^{n-1}-1\big]$
  \item $b_n=\frac12+\frac12\cdot3\big[(\frac32)^{n-1}-1\big]=\frac12+\frac32\big(\frac32\big)^{n-1}-\frac32 = \frac32\big(\frac32\big)^{n-1}-1$
  \item $a_n=2^n b_n = 2^n\big[\frac32\big(\frac32\big)^{n-1}-1\big]=3^n-2^n$
\end{itemize}

\textbf{答案}：$a_n=3^n-2^n$

\textbf{例 26}（数列与不等式）$S_n$ 为 $\{a_n\}$ 前 $n$ 项和，$a_1=1$，$a_{n+1}=\frac{n+2}{n}S_n$。
（1）求 $\{a_n\}$ 的通项公式；
（2）设 $b_n=\frac1{a_n}$，求 $\sum_{k=1}^n b_k$ 并证 $\sum_{k=1}^n b_k < 2$。

\textbf{解析}：

\textbf{（1）}：
\begin{itemize}
  \item $a_{n+1}=\frac{n+2}{n}S_n$，$a_n=\frac{n+1}{n-1}S_{n-1}\;(n\ge2)$
  \item 相减：$a_{n+1}-a_n = \frac{n+2}{n}S_n - \frac{n+1}{n-1}S_{n-1}$
  \item 利用 $S_n=S_{n-1}+a_n$ 化简得 $(n-1)a_{n+1}=(n+1)a_n$
  \item $\frac{a_{n+1}}{a_n}=\frac{n+1}{n-1}$，累乘得 $a_n=n(n+1)\;(n\ge2)$，$a_1=1$ 也符合
  \item $\therefore a_n=n(n+1)$
\end{itemize}

\textbf{（2）}：
\begin{itemize}
  \item $b_n=\frac1{n(n+1)}=\frac1n-\frac1{n+1}$
  \item $\sum_{k=1}^n b_k = 1-\frac1{n+1} < 1$（原题 $<2$ 显然成立）
\end{itemize}

\subsection{多选题专项}

\textbf{例 27}（函数性质多选）已知函数 $f(x)=x^3-3x^2+ax+b$，则
A. $f(x)$ 一定有两个极值点
B. 若 $f(x)$ 有极值点 $x_1,x_2$，则 $x_1+x_2=2$
C. 若 $f(x)$ 在 $x=1$ 处取极值，则 $a=3$
D. 若 $f(x)$ 有三个零点，则 $a<3$

\textbf{解析}：
\begin{itemize}
  \item $f'(x)=3x^2-6x+a$，$\Delta=36-12a$
  \item A：$a<3$ 时有两个极值点，$a\ge3$ 时无极值点 $\rightarrow$ \textbf{错}
  \item B：两根和 $x_1+x_2=2$ \textbf{对}
  \item C：$f'(1)=0\Rightarrow3-6+a=0\Rightarrow a=3$ \textbf{对}
  \item D：三个零点需 $\Delta>0$ 且 $f(x_1)f(x_2)<0$，仅 $a<3$ 不够 \textbf{错}
\end{itemize}

\textbf{答案}：BC

\textbf{例 28}（抽象函数多选）$f(x)$ 定义域为 $\mathbb{R}$，$f(x+y)+f(x-y)=2f(x)f(y)$，$f(0)\neq0$，则
A. $f(0)=1$ \quad B. $f(x)$ 为偶函数 \quad C. $f(x)$ 有周期 \quad D. 若 $f(1)=\frac12$，则 $f(2)=-\frac12$

\textbf{解析}：
\begin{itemize}
  \item 令 $x=y=0$：$2f(0)=2[f(0)]^2$，$f(0)\neq0\Rightarrow f(0)=1$ \textbf{A对}
  \item 令 $x=0$：$f(y)+f(-y)=2f(0)f(y)=2f(y)\Rightarrow f(-y)=f(y)$，偶函数 \textbf{B对}
  \item 令 $y=x$：$f(2x)+f(0)=2[f(x)]^2\Rightarrow f(2x)=2[f(x)]^2-1$
    $f(1)=\frac12\Rightarrow f(2)=2\cdot\frac14-1=-\frac12$ \textbf{D对}
  \item 周期需额外条件，C 不确定
\end{itemize}

\textbf{答案}：ABD

\subsection{新定义题型}

\textbf{例 29}（新定义——取整函数）设 $[x]$ 表示不超过 $x$ 的最大整数。已知 $a_n=[\sqrt{n}]$，$b_n=[\log_2 n]$。
（1）求 $S_{100}=\sum_{n=1}^{100} a_n$；
（2）求 $T_{100}=\sum_{n=1}^{100} b_n$。

\textbf{解析}：

\textbf{（1）}：
\begin{itemize}
  \item $[\sqrt{n}]=k$ 当且仅当 $k^2\le n < (k+1)^2$，这样的 $n$ 有 $2k+1$ 个
  \item $\sqrt{100}=10$，$k=1$ 到 $9$ 各 $2k+1$ 项，$k=10$ 有 $100-99=1$ 项
  \item $S_{100}=\sum_{k=1}^{9}k(2k+1)+10\times1 = \sum(2k^2+k)+10$
  \item $=2\cdot\frac{9\cdot10\cdot19}{6}+\frac{9\cdot10}{2}+10 = 570+45+10=625$
\end{itemize}

\textbf{（2）}：
\begin{itemize}
  \item $[\log_2 n]=k$ 当且仅当 $2^k\le n < 2^{k+1}$
  \item $k=0$：$1\le n<2$，1 项；$k=1$：2--3，2 项；$k=2$：4--7，4 项；$\cdots$
  \item $T_{100}=\sum_{k=0}^{6}k\cdot2^k + 7\cdot(100-2^7+1)$
  \item $=0+1\cdot2+2\cdot4+3\cdot8+4\cdot16+5\cdot32+6\cdot64 + 7\cdot37$
  \item $=2+8+24+64+160+384+259=901$
\end{itemize}

\textbf{答案}：（1）625；（2）901

\subsection{真题代练核心总结}

\begin{center}
\begin{tabular}{c|c}
\textbf{题型} & \textbf{核心方法} \\ \hline
选填压轴 & 特值法、数形结合、排除法、选项结构分析 \\
三角解答 & 边角互化、辅助角公式、正弦/余弦定理 \\
数列解答 & 构造法、累加累乘、裂项相消、错位相减 \\
立几解答 & 向量法建系、法向量、体积法辅助 \\
概统解答 & 分布列、全概率、二项分布、决策分析 \\
解几解答 & 韦达定理、设而不求、几何条件代数化 \\
导数解答 & 含参讨论、构造函数、切线放缩、隐零点 \\
多选 & 宁少勿滥、逐项验证、举反例排除 \\
新定义 & 理解定义、转化为已知模型
\end{tabular}
\end{center}

\textbf{考场策略}：
\begin{itemize}
  \item 选填卡壳超3分钟先跳过
  \item 大题第一问必须拿（通常是送分的基础计算）
  \item 解析几何重在联立+韦达，不要因计算复杂放弃
  \item 导数题求导正确就得 2--4 分，分类讨论再得 4--6 分
  \item 多选不确定的选项坚决不选（2 分 vs 0 分的取舍）
\end{itemize}
